\begin{figure*}[t]
\centering
\includegraphics[width=0.96\textwidth]{fig_rewrite_sl.pdf}
\caption{\textbf{The null replicates on a second benchmark, a second checkpoint, and a second metric family.} The SafeLIBERO counterpart of Fig.~\ref{fig:rewrite}, evaluated on \texttt{pi05\_libero\_finetuned} under the LeRobot stack. Each column holds one suite; filled circles with dashed baselines are level~I and open diamonds with dotted baselines are level~II, and each band is the 95\% Wilson interval of the unmodified benchmark instruction at that level. The in-house rewriter contributes three points per cell, one per prompt variant, and each API rewriter contributes one at \textit{v2}, giving 55 complete cells at $n{=}200$ each; the grey cross marks the single cell that did not run. \emph{Top:} collision rate, the complement of the benchmark's own collision-avoidance metric. Five of the 55 cells reach nominal significance under a Fisher exact test, all of them in Object~L2 and all in the same direction, with the largest being 83.5 to 73.0. None approaches the Bonferroni threshold of $p{=}0.0009$ that 55 comparisons require, and the smallest $p$ anywhere in the arm is 0.015. Object~L2 is therefore the one place where rewriting appears to do something, and it does not generalize to the other seven suite-levels. \emph{Bottom:} task success, plotted instead of safe success because the collision floor here is high enough that safe success is numerically the complement of the panel above it, differing by a median of 1.0 point. Task success does move: it falls below its baseline in 44 of the 55 cells, as far as 33.0 to 15.5 on Object~L1. Rewriting the instruction leaves the safety behavior statistically untouched while costing task performance. The vertical axes are windowed on each panel's own data rather than spanning zero to one hundred, because every collision baseline sits between 82 and 98\%.}
\label{fig:rewrite-sl}
\end{figure*}
